% ============================================================
%  chapters/02_background.tex
% ============================================================

\section{Education in Nigeria}
\label{sec:background}

The history of modern education in Nigeria began in the mid-nineteenth
century, when Christian missionaries introduced the ``Western'' style of
education to the region \citep{nigeria_education_factsheet}. These new
methods of schooling, along with Western bureaucratic and social norms,
were quickly adopted by those living in the southern region of the
country---notably in the coastal towns of Lagos, Calabar, and Port
Harcourt---producing some of the first African professionals, civil
servants, lawyers, and doctors. In contrast, the predominantly Muslim
North was far less welcoming of these new methods and ideas. Northern
communities rejected the Christian missionaries and the idea of Western
education out of concern that it would interfere with their faith, and
they sought to create an alternative that emphasized scholarship in
Islamic studies. The result was the creation of \emph{Islamiyya}
schools---established and operated by Muslim scholars with instruction
grounded in a core Muslim legal and religious curriculum
\citep{umar2003}.

Around the time of independence and in the years following, the Nigerian
state took an active role in expanding access to education for all
citizens, recognizing the importance of an educated population to the
country's development goals. The first generation of Nigerian universities
were established between 1948 and 1965 to meet the growing demand for an
educated citizenry. These institutions were fully funded by the federal
government, which held that it was in the national interest to ensure
that all students had equal opportunity to attend school regardless of
socioeconomic status \citep{babsfafunwa1992}. In time, as demand for
education continued to grow, more schools at all levels---primary,
secondary, and tertiary---were built across the country.

The 1970s oil boom played a major role in expanding the education sector,
given Nigeria's position as a major oil producer. Following the Nigerian-
Biafran Civil War of 1967--1970, and in response to growing inequality in
educational attainment and economic outcomes between the North and the
South and across gender lines, the federal government devised the
\textbf{Universal Primary Education (UPE) scheme} in 1976
\citep{csapo1983}. The UPE followed a 6-3-3-4 system: six years of
primary school, three years of junior secondary school, three years of
senior secondary school, and four years of tertiary education. The scheme
aimed to ensure universal primary school completion by making primary
education free and compulsory.

With these provisions, primary school enrollment grew well beyond what
policymakers had initially projected---a problem compounded by the fact
that initial enrollment estimates were based on dubious population data
from the 1963 census \citep{csapo1983}. The resulting oversubscription
led to a deterioration of school resources and a decline in the quality
of the UPE system. The number of schools was not keeping pace with the
growing population, there was a chronic shortage of qualified teachers,
and those teachers who were employed were not regularly paid
\citep{babsfafunwa1992}. These challenges were only exacerbated when oil
revenues---and thus government funding for education---declined in the
1980s. Mismanagement and corruption meant that only a fraction of the
budgeted funds were actually directed toward educational projects. By
1985, the federal government had effectively ceased funding primary
education \citep{umar2003}. The continued deterioration of the education
sector into the late 1990s meant that graduates were not adequately
equipped with the skills they were meant to acquire in school, contributing
to massive unemployment across the country.

As a remedy, the Federal Government introduced the \textbf{Universal
Basic Education (UBE) scheme} in 1999, again with the aim of providing
greater access to quality basic education, and in line with the United
Nations Millennium Development Goals. The UBE scheme set forth the
following objectives:

\begin{itemize}
    \item Providing uninterrupted access to nine years of formal education
    \item Reducing school dropout
    \item Making primary education free and compulsory
    \item Ensuring the acquisition of literacy, numeracy, and life skills
    \item Instilling a value of education in the populace \citep{ubec_online}
\end{itemize}

The Universal Basic Education Commission (UBEC) adopted as its philosophy
the ``integration of the individual into sound and effective citizenry,
and to provide equal education and opportunities for all.'' The new
scheme follows a 9-3-4 system.\footnote{The 9-3-4 system consists of nine
years of free and compulsory education---six years of primary school and
three years of junior secondary school---followed by three years of senior
secondary school and four years of tertiary education. The UBE promotes
automatic promotion through primary grades while ensuring quality through
continuous assessment.} Following nine years of free and compulsory
primary and junior secondary school, students sit the National Common
Entrance Examination (NCEE) to determine eligibility for admission to
senior secondary school.\footnote{Each state has approximately two federal
government colleges funded and managed by the federal government through
the Ministry of Education. Admission is merit-based (NCEE results) with
nominal tuition of approximately 16,000 naira. State-owned secondary
schools, funded by individual states, are generally of lower quality.
Private schools charge 160,000--320,000 naira annually but offer smaller
classes, modern equipment, and more highly qualified teachers
\citep{nigeria_education_profile}.} Following senior secondary school,
students sit the Secondary School Certificate Examination and may apply
to tertiary institutions.

\bigskip
\noindent Despite these policy efforts, deep regional and gender disparities in
educational attainment persist. Northern Nigeria continues to lag behind
the South on virtually every educational indicator, a pattern that the
Boko Haram insurgency---which began in earnest in 2009---has dramatically
intensified, as discussed in subsequent sections.
